% Bagian: second-order-logic
% Bab: sol-and-sets
% Subbagian: cardinalities

\documentclass[../../../include/open-logic-section]{subfiles}

\begin{document}

\olfileid[id]{sol}{set}{crd}

\olsection{Kardinalitas Himpunan}

\begin{explain}
Sebagaimana kita dapat mengungkapkan bahwa domain berhingga atau tak hingga,
!!{enumerable} atau !!{nonenumerable}, kita dapat mendefinisikan sifat suatu
himpunan bagian dari~$\Domain{M}$ sebagai berhingga atau tak hingga,
!!{enumerable} atau !!{nonenumerable}.
\end{explain}

\begin{prop}
Formula $\fn{Inf}(X) \ident$
\begin{multline*}
\lexists[u][(\lforall[x][\lforall[y][((X(x) \land X(y) \land
      \eq[u(x)][u(y)]) \lif \eq[x][y])]] \land {}
      \\ \lforall[x][(X(x) \lif X(u(x)))] \land {}
      \\\lexists[y][(X(y) \land \lforall[x][(X(x) \lif \eq/[y][u(x)])]])]
\end{multline*}
dipenuhi relatif terhadap suatu penugasan variabel~$s$ jika dan hanya jika
$s(X)$ tak hingga.
\end{prop}

\begin{prop}
Formula $\fn{Count}(X) \ident $
\begin{multline*}
\lnot\lexists[x][X(x)] \lor {} \\
\lexists[z][\lexists[u][(X(z) \land
    \lforall[x][(X(x) \lif X(u(x)))] \land {} \\ \lforall[Y][((Y(z) \land
      \lforall[x][(Y(x) \lif Y(u(x)))]) \lif X \subseteq Y])]])
\end{multline*}
dipenuhi relatif terhadap suatu penugasan variabel~$s$ jika dan hanya jika
$s(X)$ !!{enumerable}.
\end{prop}

Kita mengetahui dari Teorema Cantor bahwa terdapat himpunan-himpunan
!!{nonenumerable}, dan bahkan terdapat tak hingga banyak tingkat ukuran tak
hingga yang berbeda. Teori himpunan mengembangkan seluruh aritmetika
ukuran himpunan dan menetapkan bilangan kardinal tak hingga kepada himpunan.
Bilangan asli berperan sebagai bilangan kardinal yang mengukur ukuran
himpunan hingga. Kardinalitas himpunan !!{denumerable} adalah kardinalitas
tak hingga pertama, yang disebut~$\aleph_0$ (``alef-nol''). Ukuran tak hingga
berikutnya ialah~$\aleph_1$. Ini adalah ukuran terkecil yang dapat dimiliki
suatu himpunan tanpa menjadi terhitung (yakni, tanpa berukuran~$\aleph_0$).
Kita dapat mendefinisikan ``$X$ berukuran $\aleph_0$'' sebagai
$\fn{Aleph}_0(X) \liff \fn{Inf}(X) \land \fn{Count}(X)$. $X$ berukuran
$\aleph_1$ jika dan hanya jika semua himpunan bagiannya berhingga atau
berukuran~$\aleph_0$, tetapi himpunan itu sendiri tidak
berukuran~$\aleph_0$. Oleh karena itu, kita dapat mengungkapkan hal ini dengan
formula
$\fn{Aleph_1}(X) \ident \fn{Inf}(X) \land
  \lforall[Y][(Y \subseteq X \lif
  (\lnot\fn{Inf}(Y) \lor \fn{Aleph}_0(Y)))] \land \lnot
\fn{Aleph}_0(X)$. Sifat berukuran $\aleph_2$ didefinisikan secara serupa,
dan seterusnya.

Ada satu ukuran yang menarik perhatian secara khusus, yakni apa yang disebut
kardinalitas kontinuum. Ukuran ini adalah ukuran $\Pow{\Nat}$, atau, secara
ekuivalen, ukuran~$\Real$. Bahwa suatu himpunan berkardinalitas kontinuum juga
dapat diungkapkan dalam logika orde kedua, tetapi memerlukan
sedikit lebih banyak upaya.

\end{document}
